\chapter{Estatística Descritiva e Momentos}
\label{cap:probabilidade}

Este capítulo apresenta as ferramentas de estatística descritiva do \hayashi{}
como expressões da teoria das probabilidades. Cada função mapeia diretamente
para um conceito teórico, tornando o código legível tanto para quem vem da
matemática quanto para quem vem do Stata ou R.

% ─────────────────────────────────────────────────────────────────────────────
\section{Valor esperado}

O valor esperado $E(X)$ de uma variável aleatória discreta é a média ponderada
de seus valores. Para dados observados, o estimador natural é a média amostral:

\[
  \hat{E}(X) = \bar{X} = \frac{1}{n}\sum_{i=1}^{n} x_i
\]

Em \hayashi{}, \lstinline{mean} aceita uma lista ou uma coluna de DataFrame:

\begin{lstlisting}[caption={E(X) sobre lista e sobre coluna}]
// sobre lista
let xs = [2.0, 4.0, 4.0, 4.0, 5.0, 5.0, 7.0, 9.0]
display mean(xs)    // 5.0

// sobre coluna de DataFrame
load "dados.csv" as df
display mean(df, renda)
\end{lstlisting}

% ─────────────────────────────────────────────────────────────────────────────
\section{Variância e desvio padrão}

A variância amostral e o desvio padrão amostral são:

\[
  \widehat{\mathrm{Var}}(X) = s^2 = \frac{1}{n-1}\sum_{i=1}^{n}(x_i - \bar{X})^2
  \qquad
  s = \sqrt{s^2}
\]

O divisor $n-1$ (correção de Bessel) produz um estimador não viesado de
$\sigma^2$. Todas as funções de \hayashi{} usam essa convenção.

\begin{lstlisting}[caption={Variância e desvio padrão}]
display variance(df, renda)   // s²
display sd(df, renda)         // s
\end{lstlisting}

\begin{nota}
  \lstinline{sd} e \lstinline{std} são aliases — ambos calculam o desvio padrão
  amostral com divisor $n-1$.
\end{nota}

Para obter ambos de uma vez, \lstinline{summarize} retorna um dicionário com
todos os momentos:

\begin{lstlisting}
let stats = summarize(df, renda)

display stats["mean"]      // média
display stats["sd"]        // desvio padrão
display stats["variance"]  // variância
display stats["min"]       // mínimo
display stats["max"]       // máximo
display stats["N"]         // número de observações
\end{lstlisting}

% ─────────────────────────────────────────────────────────────────────────────
\section{Momentos condicionais}

O valor esperado condicional $E(X \mid C)$ restringe o cálculo a um subconjunto
da amostra sem criar um DataFrame auxiliar. A opção \lstinline{if =} está
disponível em todas as funções de agregação escalar:

\begin{lstlisting}[caption={Momentos condicionais com if =}]
// E(renda | setor == "publico")
display mean(df, renda, if = setor == "publico")

// Var(renda | setor == "publico")
display variance(df, renda, if = setor == "publico")

// DP(renda | setor == "publico")
display sd(df, renda, if = setor == "publico")
\end{lstlisting}

A condição é avaliada linha a linha no DataFrame — qualquer expressão booleana
sobre colunas é válida:

\begin{lstlisting}[caption={Condições compostas}]
display mean(df, salario, if = educ >= 12 and sexo == 1)
display sd(df, renda,   if = uf == "RS" or uf == "SC")
\end{lstlisting}

\begin{dica}
  \lstinline{if =} não altera o DataFrame. Para criar uma amostra permanente
  use \lstinline{filter}; para um cálculo pontual use \lstinline{if =}.
\end{dica}

% ─────────────────────────────────────────────────────────────────────────────
\section{Covariância}

A covariância amostral entre $X$ e $Y$ mede a variação conjunta:

\[
  \widehat{\mathrm{Cov}}(X, Y)
  = \frac{1}{n-1}\sum_{i=1}^{n}(x_i - \bar{X})(y_i - \bar{Y})
\]

\begin{lstlisting}[caption={Covariância entre duas colunas}]
display cov(df, educ, salario)

// condicional: apenas trabalhadores formais
display cov(df, educ, salario, if = formal == 1)
\end{lstlisting}

\begin{nota}
  $\mathrm{Cov}(X, X) = \mathrm{Var}(X)$. Para verificar:
  \lstinline{cov(df, x, x)} devolve o mesmo que \lstinline{variance(df, x)}.
\end{nota}

% ─────────────────────────────────────────────────────────────────────────────
\section{Correlação}

\subsection{Escalar entre duas variáveis}

O coeficiente de correlação de Pearson normaliza a covariância pelo produto dos
desvios padrão:

\[
  \rho(X, Y) = \frac{\mathrm{Cov}(X,Y)}{\mathrm{DP}(X)\,\mathrm{DP}(Y)}
  \in [-1,\, 1]
\]

\lstinline{corr_pair} devolve esse escalar diretamente:

\begin{lstlisting}[caption={Correlação escalar}]
display corr_pair(df, educ, salario)

// condicional: setor privado
display corr_pair(df, educ, salario, if = setor == "privado")
\end{lstlisting}

\subsection{Matriz de correlação}

Para explorar todas as relações entre múltiplas variáveis use
\lstinline{correlate}:

\begin{lstlisting}[caption={Matriz de correlação completa}]
// todas as variáveis numéricas
correlate(df)

// subconjunto explícito
correlate(df, educ, exp, salario)

// com p-values (* p<0.1  ** p<0.05  *** p<0.01)
pwcorr(df, educ, exp, salario)
\end{lstlisting}

\begin{nota}
  \lstinline{correlate} imprime a matriz mas não retorna um valor — use
  \lstinline{corr_pair} quando precisar do escalar em um cálculo.
\end{nota}

% ─────────────────────────────────────────────────────────────────────────────
\section{Mediana e quantis}

A mediana é o valor central da distribuição ordenada, menos sensível a valores
extremos que a média.

\begin{lstlisting}[caption={Mediana}]
display median(df, renda)

// condicional
display median(df, renda, if = regiao == "Sul")
\end{lstlisting}

Quantis generalizam a mediana para qualquer ponto $p \in [0, 1]$ da
distribuição:

\begin{lstlisting}[caption={Quantis}]
display quantile(df, renda, 0.25)   // 1º quartil
display quantile(df, renda, 0.50)   // mediana
display quantile(df, renda, 0.75)   // 3º quartil
display quantile(df, renda, 0.90)   // percentil 90
\end{lstlisting}

\lstinline{quantile} aceita \lstinline{if =}:

\begin{lstlisting}
// Percentil 90 da renda entre trabalhadores com ensino superior
display quantile(df, renda, 0.90, if = educ >= 16)
\end{lstlisting}

Ambas as funções aceitam listas diretamente:

\begin{lstlisting}
let xs = [3.0, 1.0, 4.0, 1.0, 5.0, 9.0, 2.0, 6.0]
display median(xs)           // 3.5
display quantile(xs, 0.75)   // 5.25
\end{lstlisting}

% ─────────────────────────────────────────────────────────────────────────────
\section{Sumário detalhado}

\lstinline{summarize} com \lstinline{detail=true} expande o dicionário de
retorno para incluir assimetria, curtose e toda a grade de percentis:

\begin{lstlisting}[caption={summarize com detail}]
let s = summarize(df, renda, detail=true)

display s["mean"]     // média
display s["sd"]       // desvio padrão
display s["p25"]      // 1º quartil
display s["p50"]      // mediana
display s["p75"]      // 3º quartil
display s["p90"]      // percentil 90
display s["p99"]      // percentil 99
display s["skew"]     // assimetria
display s["kurt"]     // curtose
\end{lstlisting}

% ─────────────────────────────────────────────────────────────────────────────
\section{Referência rápida}

\begin{center}
\begin{tabular}{lll}
\toprule
\textbf{Notação} & \textbf{Hayashi} & \textbf{Formas} \\
\midrule
$E(X)$
  & \lstinline{mean(df, x)}
  & lista / df / \lstinline{if =} \\
$\mathrm{Var}(X)$
  & \lstinline{variance(df, x)}
  & lista / df / \lstinline{if =} \\
$\mathrm{DP}(X)$
  & \lstinline{sd(df, x)}
  & lista / df / \lstinline{if =} \\
$\mathrm{Cov}(X, Y)$
  & \lstinline{cov(df, x, y)}
  & df / \lstinline{if =} \\
$\rho(X, Y)$
  & \lstinline{corr_pair(df, x, y)}
  & df / \lstinline{if =} \\
$\mathrm{Med}(X)$
  & \lstinline{median(df, x)}
  & lista / df / \lstinline{if =} \\
$Q_p(X)$
  & \lstinline{quantile(df, x, p)}
  & lista / df / \lstinline{if =} \\
\midrule
\multicolumn{2}{l}{\lstinline{correlate(df, x, y, ...)}}
  & matriz de correlação (display) \\
\multicolumn{2}{l}{\lstinline{summarize(df, x)}}
  & dicionário completo de momentos \\
\bottomrule
\end{tabular}
\end{center}

% ─────────────────────────────────────────────────────────────────────────────
\section{Exemplo completo}

O trecho abaixo ilustra o uso conjunto dessas funções em uma análise
exploratória típica antes de estimar um modelo:

\begin{lstlisting}[caption={Análise exploratória — renda e educação}]
load "pnad.csv" as df

// visão geral
summarize(df, renda)
summarize(df, educ)
correlate(df, renda, educ, exp)

// momentos por grupo
display mean(df, renda, if = sexo == 1)      // homens
display mean(df, renda, if = sexo == 0)      // mulheres

display sd(df, renda, if = sexo == 1)
display sd(df, renda, if = sexo == 0)

// distância entre mediana e média: sinal de assimetria
display mean(df, renda)
display median(df, renda)

// associação linear
display cov(df, educ, renda)
display corr_pair(df, educ, renda)

// concentração no topo
display quantile(df, renda, 0.90)
display quantile(df, renda, 0.99)
\end{lstlisting}
